We next asked how the spread of selling depended on the number of
interacting traders. With technical traders held, we compared markets
whose grids had $L=32$, 48, 64, 96, and 128 rows and columns, adding
tested fractions near the rise in majority probability at each size and
using the same 50 starting markets at every added fraction
(Fig.~\ref{fig:finite-size}). We repeated the comparison with technical
traders active at dealer depths $h=1$ and 100. In larger grids, majority
probability rose more steeply with the fraction $f$ and at smaller
fractions. With technical traders held, the fraction at which half the
markets reached a majority within 3,000 steps, $f_{50}$, fell from 0.216
at grid size $L=32$ to 0.166 at $L=128$. At every size, both active
midpoints were within 0.007 of the held value, usually at or below it,
and at $L=128$ they were within 0.001
(Fig.~\ref{fig:finite-size}C). These midpoints changed little when a
logistic curve replaced interpolation: across all size and
technical-trader conditions, the two estimates differed by at most
0.0035, less than a tenth of the shift across grid sizes
(Supplementary Tables~\ref{tab:finite-size-scaling} and~\ref{tab:revision-fits}).
The width $f_{75}-f_{25}$, the range of fractions over which the held
majority probability rose from 0.25 to 0.75, narrowed about threefold
over the same sizes (Fig.~\ref{fig:finite-size}A). With 95\% intervals
from resampling whole starting markets in brackets, it was
0.023 [0.018, 0.027] at $L=32$ and 0.008 [0.004, 0.010] at $L=128$
(Supplementary Table~\ref{tab:finite-size-scaling}).
The susceptibility multiplies the variance of the final mean Sell
fraction between starting markets by the number of interior traders, so it could not exceed 0.25 if each trader ended in Sell independently,
because the variance of an average over $N$ independent traders is at
most $1/(4N)$.
With technical traders held, its largest value over the tested fractions
was 21.2 at grid size $L=32$ and 241.2 at $L=128$
(Fig.~\ref{fig:finite-size}B,D). These values, far above that bound,
showed that the traders within each market ended in Sell together rather
than independently. Observation
time also moved the midpoint. At grid size $L=64$, extending the held
markets threefold, from 3,000 to 9,000 steps, lowered it from 0.178 to
0.167. This shift of 0.011 was similar to the change of 0.012 from $L=64$ to
$L=128$ at 3,000 steps, so over these ranges the midpoint depended on observation time
about as much as on grid size. Longer observation lowered the midpoint at each size tested this
way, $L=32$, 64, and 128, including when the Sell fraction required to
count a run was set to 0.4 or 0.6 instead of 0.5
(Supplementary Table~\ref{tab:revision-horizons}). Thus, larger grids
sharpened the rise in majority probability and moved it to smaller
fractions. The position of the rise also moved with observation time,
so the size
comparison located a crossover for a finite observation period rather than a fraction that would hold for indefinitely long runs, and
these runs did not determine whether such a fraction exists.
